% WELCOME -- the notes' section 0.3, sec.learningPython. Budget 4 min.

\subsection{What we learn about Python}

\begin{frame}
  \frametitle{The stack}
\codehl{numpy} and \codehl{scipy} for the numerics.

\pause

\codehl{pandas} and \codehl{pyarrow} for tabular data,
and \codehl{pandera} for the schemas that say what a table is.

\pause

\codehl{scikit-learn} for the fitting,
and \codehl{matplotlib} and \codehl{plotly} for the figures.
\end{frame}

\begin{frame}
  \frametitle{Four habits, taught alongside the stack}
We reason about what an implementation must satisfy before we write it:
the invariants it preserves, the round trips it closes,
and the alternative routes through a computation that must reach the same value.

\pause

Those properties are then written as tests, before the implementation they specify.

\pause

What defines an object is written down with the object:
a type carries the schema that validates it,
and a model carries its parametrisation as a frozen object, never as a computed default.

\pause

The clear implementation comes first and is understood before anything faster is written.
An optimisation is a second implementation shown to beat the first,
and the comparison is part of the material.
\end{frame}
